% Motivation
\section{Motivation}

\begin{frame}{Motivation: Why Event Severity?}
  \begin{columns}[T]
    \begin{column}{0.56\textwidth}
      \begin{itemize}
        \item \textbf{Non-recurring congestion} (accidents, breakdowns, public
              events) is the disruptive kind, and the hardest to plan for.
        \item When an event is first reported, the control room must decide
              \textbf{how much to send} (officers, barricades, diversions)
              \key{before anyone has seen the site}.
        \item That decision is currently made on \key{experience alone}.
        \item Under- or over-deploying both carry real cost: congestion on one
              side, wasted manpower on the other.
      \end{itemize}
    \end{column}
    \begin{column}{0.42\textwidth}
      \begin{block}{Operational impact}
        \begin{itemize}
          \item[\color{astramGreen}$\checkmark$] Faster first response
          \item[\color{astramGreen}$\checkmark$] Right-sized deployment
          \item[\color{astramGreen}$\checkmark$] Consistent decisions across shifts
        \end{itemize}
      \end{block}
      \begin{exampleblock}{Dataset at a glance}
        \small
        8{,}173 Bengaluru traffic events\\
        46 raw attributes $\rightarrow$ 30 engineered features\\
        4 severity classes
      \end{exampleblock}
    \end{column}
  \end{columns}
\end{frame}

% =====================================================================
